\section{Uninformed Learning in Normal-Form Games}\label{h1:finite-games}

In this section, we focus on the uninformed setting and present our results for normal-form games, where $\calX$ and $\calY$ are standard basis sets in $\dR^{m_x}$ and $\dR^{m_y}$. In this case, $A$ is the standard utility matrix, with utilities given by $u(e_i,e_j)=A_{ij}$.
Mixed strategies correspond to vectors $p,q$ in the simplexes $\simplex^{m_x}, \simplex^{m_y}$, with expected utility $u(p,q)=\inner{p}{Aq}$.



\subsection{Results and Analysis for Tsallis-INF}\label{h2:tsallis-inf}
We consider using the Tsallis-INF algorithm of~\citet{zimmert2021tsallisinf} in this setting. 
Tsallis-INF was originally proposed for the multi-armed bandit problem and shown to be optimal for both the stochastic setting and the adversarial setting.
Since then, it has been extended and analyzed in many other settings, including self-play in games~\citep{ito2025instancedependent}.
However, in all previous studies, achieving logarithmic regret bounds requires either a stationary or a self-play environment, 
and it is unclear at all whether it still works in our problem where the environment is controlled by an arbitrary opponent.
Somewhat surprisingly, we show that it does still achieve a certain game-dependent logarithmic regret in our problem.


Specifically, Tsallis-INF is an instance of the well-known Follow-the-Regularized-Leader (FTRL) framework.
It uses the negative Tsallis entropy as the regularizer, defined as (for a parameter $\alpha \in (0,1)$):
\begin{equation}
\varphi_{\alpha}(p)=\frac{1}{\alpha} \sum_{i} \rbr{p_i^\alpha - p_i},\label{eq:tsallis}
\end{equation}
along with standard importance-weighted reward estimators and a simple time-varying learning rate $\eta_t$.
See \Cref{alg:tsallis-inf} for the complete pseudocode.

\begin{algorithm}[t]
\caption{Tsallis-INF}\label{alg:tsallis-inf}
\begin{algorithmic}[1]
\Require A Tsallis entropy parameter $\alpha\in [0, 1]$ and a sequence of learning rates $\{\eta_t\}$.
\For{$t=1,2,\dots$}
  \State Let $p!t = \argmax_{p\in \simplex^{m_x}} \cbrm[\big]{\inner{p}{\sum_{s<t} g!s} + \frac{1}{\eta_t} \varphi_\alpha(p)}.$
  \State Sample $x!t\in \calX$ from the distribution specified by $p!t$ and play it.
  \State Observe the reward $r_t$.
  \State Set the reward estimator $g!t\in \dR^{\calX}$, where
      $g!t_x = \begin{cases}
          1-(1-r_t) / p!t_x, & \text{if } x!t=x, \\
          1, & \text{otherwise.}
      \end{cases}$
\EndFor
\end{algorithmic}
\end{algorithm}

We show that Tsallis-INF achieves different game-dependent logarithmic PSMR bounds for the case when: a) the game has a strict PSNE, or b) no PSNE exists. If a strict PSNE $(x^*, y^*)$ exists, the bound depends on the following sub-optimality gap parameters:
\begin{align*}
    \DeltaR_x &= u(x^*,y^*)-u(x,y^*), \quad
    \DeltaC_y = u(x^*,y)-u(x^*,y^*).
\end{align*}
From the definition of strict PSNEs, we know that these parameters are strictly positive. We further define $\DRmin = \min_{x\neq x^*}\DeltaR_x$ and $\DCmin = \min_{y\neq y^*}\DeltaC_y$. 
We point out that these parameters also characterize the sample complexity of finding a PSNE~\citep{maiti2024nearoptimal} and the external regret for self-play~\citep{ito2025instancedependent}.

On the other hand, if PSNEs do not exist, we know that the Nash value $\vMix$ is strictly higher than the pure-strategy maximin value $\vStar$, and our bound depends on the gap between them: $\DeltaMix\defeq\vMix-\vStar>0$.
More concretely, our main result is summarized in the following theorem.

\begin{theorem}\label{thm:tsallis} 
    In every game $u(\cdot, \cdot)$, against any adaptive adversary, Tsallis-INF with $\alpha=\frac 12$ and $\eta_t = \frac{1}{2\sqrt{t}}$ achieves a regret bound of $\psmr_T = \order\rbrm[\big]{\sqrt{m_x T}}$. Additionally,

    \begin{itemize}[topsep=0pt,noitemsep]
        \item If the game has a strict PSNE, we further have $\psmr_T = \order\rbrm[\big]{\frac{1}{\DCmin} \sum_{x\neq x^*} \frac{\log T}{\DeltaR_x}}$.
        \item If the game has no PSNEs, we further have $\psmr_T = \order\rbrm[\big]{\frac{m_x }{\DeltaMix}}$.
    \end{itemize}
\end{theorem}




\begin{proof}
Since our model is a special case of adversarial bandit, Theorem~1 of \citet{zimmert2021tsallisinf} directly implies that $\psmr_T\leq \ER_T\leq 8\sqrt{m_x T}+2$ (the difference in constant factors is due to our reward range $[-1,1]$ being twice as wide as theirs).

For the instance-dependent PSMR in the no-PSNE case, we observe that $\NR_T-\psmr_T=\rbrm[\big]{\vMix-\vStar}T=\DeltaMix T$. Together with the $\ER_T$ bound, we have
\begin{align*}
    \psmr_T=\NR_T-\DeltaMix T\leq \ER_T-\DeltaMix T=2+8\sqrt{m_x T}-\DeltaMix T.
\end{align*}
The function $T\mapsto 8\sqrt{m_x T}-\DeltaMix T$ is maximized at $T=\frac{16 m_x}{(\DeltaMix)^2}$ with a value of $\frac{16 m_x}{\DeltaMix}$ (see \Cref{lem:sqrt-func}), so $\psmr_T\leq \frac{16 m_x}{\DeltaMix}+2$.

For the instance-dependent PSMR in the strict PSNE case, let $(x^*, y^*)$ be the strict PSNE and $C\in [0, T]$ be the expected number of times the adversary deviates from the NE, that is, 
$
C\defeq\EE\sbrm[\Big]{\sum_{t=1}^T \one[y!t \neq y^*]}.
$
We then have
\begin{align*}
    \ER_T(x^*)-\psmr_T
    & = \EE\sbrm[\bigg]{\sum_{t=1}^T \rbrm[\Big]{u(x^*,y!t)-\vStar}}
      = \EE\sbrm[\bigg]{\sum_{t=1}^T \rbrm[\Big]{u(x^*,y!t)-u(x^*,y^*)}}\\
    & = \EE\sbrm[\bigg]{\sum_{t=1}^T \one[y!t \neq y^*] \DeltaC_{y!t}}
      \geq C \DCmin. \yestag\label{eq:er-nr-difference}
\end{align*}

Further,
\begin{align*}
    \EE\bigg[
        \sum_{t=1}^T \anglem[\big]{\DeltaR,p!t}
        \bigg] - \ER_T(x^*)
    & = \EE\sbrm[\bigg]{\sum_{t=1}^T \DeltaR_{x!t}} 
      - \EE\sbrm[\bigg]{\sum_{t=1}^T \rbrm[\Big]{u(x^*,y!t)-u(x!t,y!t)}}
      \\
    & = \EE\bigg[
        \sum_{t=1}^T \rbrm[\Big]{{
        u(x^*,y^*)-u(x^*,y!t)+u(x!t,y!t)-u(x!t,y^*)}}
        \bigg].
\end{align*}
Each summand in the RHS is equal to $0$ when $y!t=y^*$, and bounded by $4$ otherwise, so the entire expectation is bounded as $\EE\sbrm[\big]{
        \sum_{t=1}^T \anglem[\big]{\DeltaR,p!t}
        } - \ER_T(x^*) \leq 4C$.
Next, we apply the following key external regret bound
of Tsallis-INF (see e.g.,~\citealp[Theorem 1]{ito2025instancedependent}):
\begin{align*}
    \ER_T(x^*) 
    \leq \ER_T 
    & \leq 19 \EE\sbrm[\Bigg]{
        \sum_{t=1}^T \frac{1}{\sqrt{t}}
        \sum_{x\neq x^*} \sqrt{p!t_x}
        }.
\end{align*}
This is further bounded by Cauchy-Schwarz as
\begin{align*}
    \ER_T(x^*) & \leq 19 \EE\sbrm[\Bigg]{
        \sum_{t=1}^T
        \sum_{x\neq x^*} \sqrt{\frac{1}{t\DeltaR_x} \DeltaR_x p!t_x}
    } 
    \leq 19 \EE\sbrm[\Bigg]{\sqrt{
        \rbrm[\Bigg]{\sum_{t=1}^T
        \sum_{x\neq x^*} \frac{1}{t\DeltaR_x}} \rbrm[\Bigg]{\sum_{t=1}^T
        \sum_{x\neq x^*} \DeltaR_x p!t_x}
        }
    } \\
    & \leq 19 \EE\sbrm[\Bigg]{\sqrt{
        \sum_{x\neq x^*} \frac{H_T}{\DeltaR_x}
        \sum_{t=1}^T \anglem[\big]{\DeltaR, p!t}
        }
    }
    \leq 19 \EE\sbrm[\Bigg]{\sqrt{
        \sum_{x\neq x^*} \frac{H_T}{\DeltaR_x}
        \rbrm[\big]{\ER_T(x^*)+4C}
        }
    } \\
    & \leq 361 \sum_{x\neq x^*} \frac{H_T}{\DeltaR_x} + 38 \sqrt{C \sum_{x\neq x^*} \frac{H_T}{\DeltaR_x}},
\end{align*}
where $H_T=\sum_{t=1}^T \frac{1}{t} \leq 1+ \log T$ is the sum of a harmonic series, and the last step is a standard self-bounding argument and can be shown with
elementary algebra (\Cref{lem:self-bound}).
Combining this with \eqref{eq:er-nr-difference}, we can bound $\psmr_T$ as follows:
\begin{align*}
    \psmr_T \leq \ER_T(x^*) - C \DCmin \leq 361 \sum_{x\neq x^*} \frac{H_T}{\DeltaR_x} + 38 \sqrt{C \sum_{x\neq x^*} \frac{H_T}{\DeltaR_x}} - C \DCmin.
\end{align*}
Another application of \Cref{lem:sqrt-func} finally eliminates $C$ and proves our desired bound in the strict PSNE case,\footnote{Throughout this paper, no attempts are made to optimize global constants.}
\let\jmlrQED\relax %
\begin{align*}
     \psmr_T \leq 361 \rbrm[\Big]{1+ \frac{1}{\DCmin}} \sum_{x\neq x^*} \frac{H_T}{\DeltaR_x} \leq 361 \rbrm[\Big]{1+ \frac{1}{\DCmin}} \sum_{x\neq x^*} \frac{1 + \log T}{\DeltaR_x}.\mqed
\end{align*}
\end{proof}



From the proof of \Cref{thm:tsallis}, it is clear that the $\order\rbrm[\big]{\frac{m_x}{\DeltaMix}}$ instance-dependent bound in fact holds for any algorithms with a worst-case $\order(\sqrt{m_x T})$ external regret bound.
Also, for the case when a non-strict PSNE exists, $\psmr_T=\Omega(\sqrt{T})$ is in fact unavoidable; see \Cref{rem:sqrt-t-lb} below.

There is no prior work we can directly compare our bounds to for this uninformed setting.
However, since our results obviously also apply to the easier informed setting, which~\citet{maiti2025limitations} studied, we make the following comparisons to their bound.


As mentioned, \citet{maiti2025limitations} only considered a $2\times 2$ zero-sum game $A=\rbrm[\Big]{\begin{matrix}a & b \\ c & d\end{matrix}}$ with a unique NE,
and they achieved $\NR_T = \order\rbrm[\big]{\frac{\log T}{\DeltaEntry^3}+\frac{\log^2 T}{\DeltaEntry^2}}$ for another game-dependent constant $\DeltaEntry=\min\{|a-b|,|a-c|,|b-d|,|c-d|\}$.
Besides the fact that our results hold without the uniqueness assumption and that a worst-case guarantee  $\order(\sqrt{T})$ always holds, our bounds are also better in the following sense:
\begin{itemize}[itemsep=2pt,topsep=2pt]
    \item When the unique NE is pure, $\psmr_T$ is equal to $\NR_T$, and since $\DCmin\geq \DeltaEntry$ and $\DeltaR_x\geq \DeltaEntry$, our bound $\frac{1}{\DCmin} \sum_{x\neq x^*} \frac{\log T}{\DeltaR_x}$ is at most $\frac{\log T}{\DeltaEntry^2}$, strictly improving upon theirs.
    \item When the unique NE is mixed, our $\psmr_T$ metric is weaker than their $\NR_T$, but at the same time, our bound is also strictly better since $\frac{m_x}{\DeltaMix}\leq\frac{8}{\DeltaEntry^2}$ by the fact $\DeltaMix \geq \frac{\DeltaEntry^2}{4}$ from \Cref{lem:delta-comparison}.
\end{itemize}


\subsection{Lower Bound}\label{h2:uninformed-lb}




In the case when a strict PSNE exists, \citet{ito2025instancedependent} proved that Tsallis-INF enjoys $\ER_T=\order\rbrm[\big]{(\sum_x \frac{1}{\DeltaR_x}+\sum_y \frac{1}{\DeltaC_y})\log T}$ in the self-play setting.
While not directly comparable to our setting, it naturally makes one wonder whether our bound $\order\rbrm[\big]{\frac{1}{\DCmin} \sum_{x\neq x^*} \frac{\log T}{\DeltaR_x}}$ is improvable.
Indeed, their dependency on the gap parameters is of order $\frac{1}{\DRmin}+\frac{1}{\DCmin}$, thus strictly better than ours, which is of order $\frac{1}{\DRmin\DCmin}$.
It turns out that such $\frac{1}{\DRmin\DCmin}$ dependency is unavoidable in our setting, meaning that uninformed adversarial learning in games is inherently more difficult than self-play. This is summarized in the following theorem.



\begin{theorem}\label{thm:uninformed-lb}
    For small enough $\DRmin, \DCmin$, large enough $T$, and any uninformed learner, there exist an adversary and a game $u(\cdot, \cdot)$ with $m_x=m_y=2$ and a strict PSNE, such that the learner's PSMR is at least $\Omega\rbrm[\big]{\min\cbrm[\big]{\frac{1}{\DRmin\DCmin},\sqrt{T}}}$.
\end{theorem}

In the following, we will show a simplified, informal argument for the $\Omega(\frac{1}{\DRmin\DCmin})$ lower bound, using big-$\order$ to hide constant factors. A full proof with a careful analysis of constant factors is deferred to \Cref{h1a:proof-lower-bound}. \\

\begin{proof}(sketch).
    In this sketch, we focus on the more interesting case when $\frac{1}{\DRmin\DCmin}=\order(\sqrt T)$ and show that $\psmr_T = \Omega\rbrm[\big]{\frac{1}{\DRmin\DCmin}}$.
    Consider a game matrix
    $
        A=\rbrm[\bigg]{\begin{matrix}
            0 &        \DeltaC \\ 
            -\DeltaR & K - \DeltaR 
          \end{matrix}},
    $
    parameterized by $K$, $\DeltaR$, and $\DeltaC$. We assume that all entries are in $[-1, +1]$ and $\DeltaR, \DeltaC>0$. We know that $A$ admits a strict PSNE $(x^*,y^*)=(e_1,e_1)$ with a value of $v^*=0$. The two gap variables $\DeltaR,\DeltaC$ are exactly $\DRmin$ and $\DCmin$ of $A$. We set the noise model to be binary on $\{\pm 1\}$, i.e., at round $t$, the reward comes from a two-point distribution so that $r_t\in \{-1,+1\}$ and $\EE[r_t]=u(x!t,y!t)$.
    
    The adversary is parameterized by $\eps\in (0,1)$ and $T'\in \bbN_+$, $T'\leq T$, where he plays a sample from the distribution $y!t\sim q!t=(1-\eps,\eps)$ for $t\leq T'$ and $y!t=y^*$ for $t>T'$. Then, for $t\leq T'$, we have
    $
        A q!t = \rbrm[\bigg]{\begin{matrix}
            \eps \DeltaC\\
            K \eps - \DeltaR
        \end{matrix}}.
    $

    Let $\delta=\eps\DeltaC-K \eps+\DeltaR$. In each round up to $t=T'$, the KL divergence between the reward distributions of the two actions of the learner can be bounded by $\order(\delta^2)$ (see \Cref{lem:kl-bernoulli}). Using standard arguments (see Chapter 15.1 of \citet{lattimore2020bandit}), the total KL divergence between the two arms at the end of the $T'$-th round is bounded by $\delta^2T'$. 

    We let $K=1$ and $\eps=\DeltaR-12\DeltaR\DeltaC$, and thus $\delta=\DeltaR\DeltaC(13-12\DeltaC)=\Theta(\DeltaR\DeltaC)$. By choosing some $T'\leq T$ such that $T'=\Theta\rbrm[\big]{\frac{1}{(\DeltaR\DeltaC)^2}}$ and $\delta^2 T'\leq 1$, which is possible because $\frac{1}{\DeltaR\DeltaC}=\order(\sqrt T)$, we can see that the total KL divergence is $\order(1)$.
    This means that the learner cannot distinguish between the two actions beyond a constant error probability, and that the worse action is selected at least $\frac1{12}$ of the time. In this case, the accumulated PSMR up to time $T'$ is at least
    \begin{align*}
        \psmr_{T'}&\geq T'\rbrm[\Big]{\frac{1}{12} (\DeltaR-K\eps)-\frac{11}{12} \eps\DeltaC} \yestag\label{eq:psmr-lb-main}\\
        &\geq \frac{T'}{12}\DeltaR\DeltaC=\Theta\rbrm[\big]{\frac{1}{\DeltaR\DeltaC}}.
    \end{align*}

    Finally, note that for $t\in [T'+1,T]$, the adversary plays the PSNE action and thus the PSMR can only increase, proving that $\psmr_T$ is indeed also of order $\frac{1}{\DRmin\DCmin}$.
\end{proof}

\begin{remark}\label{rem:sqrt-t-lb}
    In the construction above, we can let $\DeltaR=\DeltaC=0$, $K=-1$, and let the adversary play $(1-\eps, \eps)$ with $\eps=\frac{1}{\sqrt T}$ for all $T'=T$ rounds. We can verify that the accumulated KL divergence is $\order(1)$ as well. 
    This matrix, $\rbrm[\Big]{\,\begin{matrix}
        0 & 0 \\ 
        0 & -1
        \end{matrix}}$,
    has a non-strict PSNE, meaning that $\DRmin=\DCmin=\DeltaMix=0$, and \eqref{eq:psmr-lb-main} shows that $\psmr_T = \Omega(\sqrt{T})$.
    This shows that, beyond the two cases considered in \Cref{thm:tsallis} where we prove game-dependent logarithmic PSMR, any algorithm must suffer $\psmr_T=\Omega(\sqrt T)$.
\end{remark}


\section{Informed Learning in Normal-Form Games}\label{h1:pure-ucb}

The previous section demonstrates a gap between adversarial and self-play learning in games.
To improve performance in the adversarial bandit setting, we turn to the informed setting and grant the algorithm access to extra side information, the opponent's action $y!t$. 
Note that previous work~\citep{odonoghue2021matrix,maiti2025limitations} has implicitly assumed this feedback model.


\begin{algorithm}[t]
\caption{\PureUCB}\label{alg:pure-ucb}
\begin{algorithmic}[1]
\Require A high-probability parameter $\delta$.
\State Initialize the accumulated rewards $R!0(x,y)=0$, and the number of pulls $N!0(x,y)=0$.
\For{$t=1,2,\dots$}
  \State Define the UCB for any pair $(x,y)$ as $U!{t}(x, y)\gets 1$ if $N!{t-1}(x,y)=0$, or else\\
  {$U!{t}(x, y)\gets
    \frac{R!{t-1}(x,y)}{N!{t-1}(x,y)}+\sqrt{\frac{4\log (1/\delta)+2\log (1+N!{t-1}(x,y))}{N!{t-1}(x,y)}}
  $} if $N!{t-1}(x,y)>0$. \label{code:ucb-confidence}
  \State Let $x!t=\argmax_{x\in \calX}\min_{y\in \calY} U!{t}(x,y)$ with ties broken arbitrarily.\label{code:ucb-maximin}
  \State Play $x!t$. Observe the adversarial action $y!t$ and the reward $r_t$.
  \State Update $N!t(x!t,y!t)\gets N!{t-1}(x!t,y!t)+1$. \label{code:ucb-n}
  \State Update $R!t(x!t,y!t)\gets R!{t-1}(x!t,y!t)+r_t$. \label{code:ucb-sum}
  \State All other entries of $N!t$ and $R!t$ are the same as $N!{t-1}$ and $R!{t-1}$.
\EndFor
\end{algorithmic}
\end{algorithm}

With such extra information, it becomes natural to directly build estimates on the game matrix $A$ and apply the optimism principle as in the standard UCB algorithm~\citep{auer2002using}.
This leads to our proposed algorithm \PureUCB{} (\Cref{alg:pure-ucb}). 
Specifically, at the beginning of each round $t$, it constructs an estimate $U!t(x,y)$ for every action pair $(x,y)$, representing an upper confidence bound of the utility $u(x,y)$. %
The algorithm then simply plays the optimistic pure maximin strategy $\argmax_{x\in \calX}\min_{y\in \calY} U!{t}(x,y)$ in round $t$, that is, the action that maximizes the worst-case optimistic reward.

We point out that  \PureUCB{} is very close to Algorithm~1 of \citet{odonoghue2021matrix},
with the only difference being that their algorithm plays the optimistic \textit{mixed} maximin strategy $\argmax_{p\in \simplex^{m_x}}\min_{y\in \calY} U!{t}(p,y)$ instead.
While they only prove that their algorithm achieves
$\NR=\tilde{\order}(\sqrt{m_x m_y T})$ (which is in fact even worse than using a standard adversarial bandit algorithm), 
we prove the following game-dependent logarithmic PSMR for our slightly different algorithm \PureUCB{}, using yet another sub-optimality gap parameter: $\Delta_{xy}=\vStar-u(x,y)$ for all $(x,y)\in \calX\times \calY$.
(Note that $\Delta_{xy}$ can be negative.)

\begin{theorem}\label{thm:pureucb}
    For any game $u$, \PureUCB{} with $\delta=\frac{1}{T}$ simultaneously achieves an instance-dependent regret of $\psmr_T=\order\rbrm[\big]{\sum_{x, y : \Delta_{xy} > 0} \rbrm[\big]{\Delta_{xy}+\frac{1}{\Delta_{xy}}\log T}}$ and a worst-case bound of $\psmr_T=\order\rbrm[\big]{\sqrt{m_x m_y T \log T}}$, against any adaptive adversary.
\end{theorem}

The proof mostly extends standard UCB proofs and argues that each $(x,y)$ pair can be encountered at most  $\tilde{\order}\rbrm{1/\Delta_{xy}^2}$ if its gap $\Delta_{xy}>0$; see \Cref{h1a:proof-pureucb} for a rigorous proof.
Note that unlike \Cref{thm:tsallis} where different cases were discussed, the results of \Cref{thm:pureucb} hold always (regardless whether a PSNE exists or not).
The instance-dependent bounds in these two theorems are generally incomparable (setting aside the fact that their feedback models are also different):
\begin{itemize}[itemsep=2pt,topsep=2pt]
    \item There are games where Tsallis-INF achieves a much better bound. For example, consider $A=\rbrm[\Big]{\begin{matrix}
        0 & 1 \\ 
        -1 & -\eps
        \end{matrix}}$ for a small $\eps > 0$,
        in which case Tsallis-INF achieves $\psmr_T=\order(\log T)$ since $\DRmin=\DCmin=1$, 
        while \PureUCB{} achieves $\order(\frac{\log T}{\eps})$ since $\Delta_{2,2} = \eps$.
    %
     
     \item Conversely, there are also games where \PureUCB{} achieves a much better bound. For example, consider $A=\rbrm[\Big]{\begin{matrix}
    0 & \eps \\ 
    -1 & 1
    \end{matrix}\,}$ for a small $\eps > 0$,
    in which case Tsallis-INF achieves $\psmr_T=\order(\frac{\log T}{\eps})$ since $\DRmin=\eps$ and $\DCmin=1$, 
    while \PureUCB{} achieves $\order(\log T)$ since the only non-negative gap is $\Delta_{2,1} = 1$.
\end{itemize}

Similarly, using the same examples, we see that \Cref{thm:pureucb} and the result of~\citet{maiti2025limitations} are also directly incomparable:
in the first example, we have $\DeltaEntry=1-\eps$ and thus \citet{maiti2025limitations} achieve $\psmr_T=\NR_T=\order(\log^3 T)$, better than \PureUCB{}'s bound $\order(\frac{\log T}{\eps})$;
in the second example, we have $\DeltaEntry=\eps$ and thus \citet{maiti2025limitations} achieve $\psmr_T=\NR_T=\order\rbrm[\big]{\frac{\log T}{\eps^3}+\frac{\log^2 T}{\eps^2}}$, much worse than \PureUCB's bound $\order(\log T)$. %

 



